\problemname{The Ferris Wheel}
\illustration{0.3}{image}{Singapore Flyer \newline © CEphoto, Uwe Aranas}%
%
After following a very well-optimized route through the airport, the Swedish team has finally arrived at IOI in Singapore.
During the first excursion, all $N$ teams at IOI get to ride the giant Ferris wheel Singapore Flyer.
The Ferris wheel has $M$ gondolas and it also takes $M$ minutes for the wheel to complete one revolution
(so it takes 1 minute for each gondola to move one step).

Some teams seem to be more interested in riding the Ferris wheel than others, and therefore each team gets to decide exactly how many revolutions they want to ride.
It would be boring for the participants if they had to get off and then back on before completing all their revolutions.
It has therefore been decided that once a team has sat down in a gondola, that team may remain in the gondola until they have completed all their revolutions.
This means that if the wheel turns so that a gondola reaches the entrance, but a team is already sitting in that gondola and wants to continue riding,
the next team cannot enter that gondola. This team must then wait for an empty gondola or a gondola with a team that is getting off.

The teams are very efficient at getting in and out of the gondolas, so it takes no extra time to switch teams in the bottom gondola.

All teams are currently queuing in front of the Ferris wheel, and the Swedish team wonders how long it will take before everyone has had their ride.

\section*{Input}
The first line contains the integers $N$ and $M$ ($1 \leq N,M \leq 2 \cdot 10^5$),
the number of teams and the number of gondolas in the Ferris wheel.

The second line contains $N$ integers $T_1, \dots , T_N$ ($1 \leq T_i \leq 10^9$), where $T_i$ is the number of revolutions team
number $i$ wants to ride. The teams are ordered by their place in the queue, where $T_1$ is the first team in the queue.

\section*{Output}
Print a line with an integer, the number of minutes it will take for all teams to ride.

\section*{Scoring}
Your solution will be tested on a set of test groups.
To earn points for a group, you must pass all test cases in that group.

\noindent
\begin{tabular}{| l | l | p{12cm} |}
  \hline
  \textbf{Group} & \textbf{Points} & \textbf{Constraints} \\ \hline
    $1$   & $20$     & $1 \leq N, M, T_i \leq 100$ \\ \hline
    $2$   & $30$     & $1 \leq N, M, T_i \leq 1000$ \\ \hline
    $3$   & $25$     & $1 \leq N, M \leq 1000$ \\ \hline
    $4$   & $25$     & No additional constraints. \\ \hline
\end{tabular}


\section*{Explanation of Samples}

\begin{figure}[h]
\includegraphics[width=0.8\textwidth]{sample1}
\caption{Sample 1}
\end{figure}

In sample $1$ there are $4$ teams and $3$ gondolas.
The teams, which in the picture are Sweden, Norway, Finland and Denmark, want to ride $2$, $2$, $1$ and $1$ revolutions respectively.
Note that the Danish team cannot enter the Ferris wheel at $t=3$ or $t=4$
because the Swedish / Norwegian team is already sitting in the bottom gondola and in both cases wants to ride one more revolution.
